%% sections/03-sample-covariance-wishart.tex
%% Section 3: Sample Covariance Matrices (~11 pages)
%% Source notes:
%%   - v5 LaTeX notes: sample covariance, Wishart, fixed-dim CLT, 2D Wishart,
%%     GOE, delta-method, Marchenko-Pastur, Tracy-Widom heuristic.
%%   - Handwritten Note 2.pdf, Note 3.pdf, Note 4.pdf.
%%   - Gustavo Notes (April 17, April 27, April 30, May 14).
%%   - Yao, Zheng, Bai (2015) Chapter 3 (Wishart) and Chapter 2 (MP).
%%   - Joseph Genzer (2025) Section 3.

\section{Sample Covariance Matrices}\label{sec:scm}

\todo{Opening paragraph: the sample covariance matrix \(\Sigmahat_n\) is the
central object of multivariate statistics. In low dimensions (\(p\) fixed,
\(n \to \infty\)) it converges almost surely to the population covariance
\(\Sigma\) and its eigenvalues converge by continuity. In high dimensions
(\(p, n \to \infty\) with \(p/n \to y > 0\)) this convergence fails
dramatically: the empirical spectrum spreads out into a non-trivial bulk
density. This section develops both regimes, building from the foundations
laid in Section~\ref{sec:foundations}.}

\subsection{Population and Sample Covariance, Spectral Theorem, Positive Semidefiniteness}\label{subsec:scm-foundations}

\todo{Content to fill in (from v5 + Joseph Section 3 intro):
- Definition: sample covariance matrix \(\Sigmahat_n = \frac{1}{n}
  \sum_{i=1}^n \vect{X}_i \vect{X}_i\transpose\) (mean-zero case) and the
  unbiased \(\frac{1}{n-1}\) variant.
- Matrix form: \(\Sigmahat_n = \frac{1}{n} \mat X \mat X\transpose\) for
  \(\mat X \in \R^{p \times n}\).
- Spectral theorem for symmetric matrices: \(A = Q \Lambda Q\transpose\).
- PSD characterization: \(A\) is PSD iff all eigenvalues are nonneg (proof).
- Use the diagonalization to compute quadratic forms; foreshadow why
  eigenvalues of \(\Sigmahat_n\) are the meaningful quantities.}

\subsection{Low Dimensions: Wishart, Eigenvalue Repulsion, and Fluctuations}\label{subsec:scm-lowdim}

\todo{Content to fill in (from v5 + Joseph Section 3.1):

(a) Almost sure convergence of the sample covariance.
- Theorem: if \(\vect X_i\) iid with \(\E \vect X_1 = 0\) and \(\Cov(\vect X_1)
  = \Sigma\), then \(\Sigmahat_n \convas \Sigma\) entrywise by SLLN.
- Corollary: \(\lambda_i(\Sigmahat_n) \convas \lambda_i(\Sigma)\) by
  continuity of eigenvalues (proof sketch via continuity of the
  characteristic polynomial and continuity of its roots).
Cite \citet[Prop.~3.1, 3.3]{Genzer2025}.

(b) Wishart distribution.
- Definition: \(\mat S = \sum_{i=1}^n \vect X_i \vect X_i\transpose \sim
  \Wishart{p}{n, \Sigma}\) for iid Gaussian \(\vect X_i \sim \Normal{0,
  \Sigma}\).
- Wishart density formula.
- Theorem: for \(p = 1\), \(\Wishart{1}{n, 1}\) collapses to \(\chi_n^2\)
  (proof by plugging in).

(c) 2D Wishart and eigenvalue repulsion.
- Joint density of ordered eigenvalues of \(\Wishart{2}{n, I_2}\) contains
  the factor \(\lambda_2 - \lambda_1\), producing eigenvalue repulsion.
- Simulation: scatter of \((\lambda_1, \lambda_2)\) and histogram of
  \(\lambda_2 - \lambda_1\) (notebook 02).
Include figure: 2D Wishart eigenvalue scatter.

(d) Joint CLT for entries of \(\Sigmahat_n\) when \(\Sigma = I_2\).
- Statement and proof via multivariate CLT (variances 2, 2, 1 for diagonal
  and off-diagonal entries).
- Consequence: matrix-valued CLT \(\sqrt{n}(\Sigmahat_n - I_2) \convd G\)
  where \(G\) has GOE structure.

(e) Top eigenvalue fluctuation: repeated vs.\ distinct cases.
- Repeated (\(\Sigma = I_2\)): \(\sqrt{n}(\lambda_2(\Sigmahat_n) - 1) \convd
  \lambda_2(G)\), \emph{not} Gaussian.
- Distinct (\(\Sigma = \diag(1, 2)\)): delta-method gives Gaussian
  fluctuations with variances 2 and 8.
- Simulation: Q-Q plots for both cases (notebook 03).
Include figures: Q-Q repeated, Q-Q distinct.

(f) Eigenvector convergence.
- Distinct case: \(\widehat{\vect u}_\ell \convas \vect e_\ell\) with sign
  convention.
- Repeated case: eigenvectors fail to converge; explicit counterexample
  with rotating \(O_n\). This previews why spectral separation is essential
  for eigenvector recovery, a theme that returns in Section~\ref{sec:spiked}.}

\subsection{High Dimensions: Marchenko-Pastur and the Tracy-Widom Edge}\label{subsec:scm-highdim}

\todo{Content to fill in (from v5 high-dim section + Joseph Section 3.2):

(a) The high-dimensional regime.
- Setup: \(p, n \to \infty\) with \(p/n \to y \in (0, \infty)\).
- Empirical spectral distribution \(F_{\Sigmahat_n}(\lambda) = \frac{1}{p}
  \#\{i : \lambda_i(\Sigmahat_n) \le \lambda\}\).

(b) Marchenko-Pastur law.
- Theorem (Marchenko-Pastur 1967): for null \(\Sigma = I_p\) and \(0 < y < 1\),
  \(F_{\Sigmahat_n} \convd \mathrm{MP}_y\) with support
  \([(1 - \sqrt{y})^2, (1 + \sqrt{y})^2]\) and density
  \(f_{\mathrm{MP}}(\lambda) = \frac{1}{2\pi y \lambda}
   \sqrt{(\lambda_+ - \lambda)(\lambda - \lambda_-)}\).
- For \(y \ge 1\), an atom at zero of mass \(1 - 1/y\).
- Convergence of \(\lambda_{\max}(\Sigmahat_n) \convas (1 + \sqrt{y})^2\).
- Simulation: histogram against MP density at \(y = 1/2\) (notebook 04).
Include figure: MP density vs.\ simulation.

(c) Tracy-Widom edge heuristic.
- Goal: explain why edge fluctuations live on scale \(p^{-2/3}\), not
  \(n^{-1/2}\).
- Heuristic: square-root density near the edge,
  \(f_{\mathrm{MP}}(\lambda) \sim C(\lambda_+ - \lambda)^{1/2}\) as
  \(\lambda \uparrow \lambda_+\), combined with the requirement that the
  mass to the right of \(\lambda_{\max}\) be of order \(1/p\), yields
  \(\lambda_+ - \lambda_{\max} \sim p^{-2/3}\).
- Statement (informal): \(p^{2/3}(\lambda_{\max} - \lambda_+) \convd
  \mathrm{TW}_1\) after model-dependent constants.
- Remark on large-deviation heuristic confirming the same scaling from
  left and right tails.
Cite \citet{TracyWidom1996,Johnstone2001}.}

\clearpage
